% Gerado por regressao/06_tabelas_latex.R (projeto Modelagem). Não editar à mão.
\begin{sidewaystable}
\caption{Robustez R5 — especificação de Sakurai (sem pré-eleitoral e sem pandemia, erro agrupado)}\label{tab:robustez_R5}
\centering
\footnotesize
\setlength{\tabcolsep}{3pt}
\begin{tabular}{@{}>{\raggedright\arraybackslash}p{\dimexpr\linewidth-594pt\relax}>{\centering\arraybackslash}p{60pt}>{\centering\arraybackslash}p{60pt}>{\centering\arraybackslash}p{60pt}>{\centering\arraybackslash}p{60pt}>{\centering\arraybackslash}p{60pt}>{\centering\arraybackslash}p{60pt}>{\centering\arraybackslash}p{60pt}>{\centering\arraybackslash}p{60pt}>{\centering\arraybackslash}p{60pt}@{}}
\hline
\textbf{Variável} & \textbf{\hspace{0pt}Despesa total} & \textbf{\hspace{0pt}Saúde e Saneamento} & \textbf{\hspace{0pt}Educação e Cultura} & \textbf{\hspace{0pt}Habitação e Urbanismo} & \textbf{\hspace{0pt}Assistência Social} & \textbf{\hspace{0pt}Transporte} & \textbf{\hspace{0pt}Administração} & \textbf{\hspace{0pt}Agricultura} & \textbf{\hspace{0pt}Comunicações} \\ \hline
\multicolumn{10}{@{}l}{\textit{Modelo A (ciclo eleitoral)}} \\
Ano eleitoral & 0,0432\textsuperscript{***} & 0,0689\textsuperscript{***} & $-$0,0257\textsuperscript{***} & 0,2561\textsuperscript{***} & 0,0604\textsuperscript{***} & 0,1909\textsuperscript{***} & 0,0152\textsuperscript{***} & 0,0423\textsuperscript{***} & 0,0227 \\
 & (0,0016) & (0,0020) & (0,0019) & (0,0054) & (0,0025) & (0,0089) & (0,0024) & (0,0059) & (0,0144) \\
Prefeito na coalizão & 0,0022\textsuperscript{**} & 0,0011 & 0,0013 & 0,0113\textsuperscript{*} & $-$0,0014 & 0,0093 & 0,0059\textsuperscript{**} & 0,0075 & $-$0,0053 \\
do governador & (0,0010) & (0,0018) & (0,0017) & (0,0063) & (0,0029) & (0,0100) & (0,0026) & (0,0075) & (0,0174) \\
Prefeito na coalizão & 0,0084\textsuperscript{***} & $-$0,0049\textsuperscript{**} & 0,0309\textsuperscript{***} & $-$0,0231\textsuperscript{***} & 0,0065\textsuperscript{*} & 0,0358\textsuperscript{***} & 0,0034 & 0,0091 & 0,0053 \\
do presidente & (0,0012) & (0,0022) & (0,0020) & (0,0077) & (0,0036) & (0,0123) & (0,0033) & (0,0090) & (0,0217) \\
\hline
Observações & 70.669 & 70.526 & 70.550 & 69.574 & 70.484 & 54.422 & 70.553 & 64.464 & 12.360 \\
Municípios & 5.568 & 5.568 & 5.568 & 5.564 & 5.568 & 5.205 & 5.568 & 5.436 & 1.896 \\
R\textsuperscript{2} & 0,850 & 0,676 & 0,659 & 0,262 & 0,377 & 0,136 & 0,305 & 0,132 & 0,029 \\
\hline
\end{tabular}
\fonte{Elaboração própria.}
\nota{Despesa e receitas \textit{per capita}, em R\$ de 2025, em logaritmo (\autoref{sec:especificacao}). O coeficiente de uma \textit{dummy} se lê como variação percentual aproximada (exata: 100 $\times$ (exp(b) $-$ 1)). Só modelo A: efeito fixo de município, sem ano pré-eleitoral e sem pandemia. Erros-padrão agrupados por município entre parênteses. \textsuperscript{***}p$<$0,01; \textsuperscript{**}p$<$0,05; \textsuperscript{*}p$<$0,1.}
\end{sidewaystable}
